\documentclass[11pt]{article}

\usepackage{fontspec}
\usepackage{polyglossia}
\setmainlanguage{french}
\setmainfont{Latin Modern Roman}

\usepackage[margin=2.6cm]{geometry}
\usepackage{booktabs}
\usepackage{array}
\usepackage{amsmath}
\usepackage{amssymb}
\usepackage{enumitem}
\usepackage{titlesec}
\usepackage[hidelinks]{hyperref}
\usepackage{tcolorbox}
\tcbuselibrary{breakable}
\usepackage{tikz}
\usetikzlibrary{arrows.meta, positioning, calc, decorations.pathreplacing}

\titleformat{\section}{\Large\bfseries}{\thesection}{0.7em}{}
\titleformat{\subsection}{\large\bfseries}{\thesubsection}{0.6em}{}

\newtcolorbox{remarque}[1][Remarque]{breakable, colback=black!3, colframe=black!40,
  boxrule=0.4pt, arc=1pt, title=#1, fonttitle=\bfseries, coltitle=black}

\setlength{\parskip}{0.5em}
\setlength{\parindent}{0pt}

\title{Résultat A, au complet\\[0.3em]
\large Récupération de la couche de sortie d'un réseau profond, en hard-label,
indépendamment de l'activation des couches cachées}
\author{Tidiane Diallo}
\date{5 novembre 2026}

\begin{document}
\maketitle

\begin{remarque}[Ce que couvre ce document]
Le document envoyé le 18 octobre résumait deux résultats, A et B, sur une
page chacun. Pr.~Ciss a demandé le détail complet du résultat A. C'est ce que
ce document fournit : le problème, les trois hypothèses sur lesquelles repose
la méthode, le déroulement pas à pas sur un cas réel, les résultats mesurés
sur les quatre activations testées, et les limites. Le résultat B (extraction
hard-label de neurones non ReLU, encore un prototype) n'est pas repris ici.

Tous les chiffres de ce document proviennent de l'exécution réelle du script
\texttt{attaque\_couche\_sortie.py}, réexécutée pour la rédaction de ce
document. Les journaux complets des quatre exécutions sont dans le dossier
\texttt{logs/} déjà transmis avec le document du 18 octobre.
\end{remarque}

\tableofcontents

\section{Le problème}

Carlini et al. (EUROCRYPT 2025) savent récupérer toutes les couches d'un
réseau ReLU en hard-label, sauf la dernière. La couche de sortie produit des
logits bruts, sans passer par une activation : il n'y a donc pas de coude à
détecter, et la technique employée pour les couches cachées ne s'y applique
pas.

Canales-Martínez et Santos (2025), mon article de départ, comblent ce trou,
mais seulement dans un cadre où le reste du réseau est en ReLU. La question
posée ici est simple à énoncer : est-ce que la récupération de la couche de
sortie dépend réellement de ReLU, ou seulement de la règle de décision
finale, qui est la même quelle que soit l'activation cachée ?

La réponse mesurée dans ce document est non : la méthode fonctionne à
l'identique sur ReLU, GELU, SiLU et tanh, avec les mêmes garanties
numériques. Le reste du document explique pourquoi, puis le prouve par
l'exécution.

\section{Les trois hypothèses}

Soit $y \in \mathbb{R}^{d_r}$ la sortie de la dernière couche cachée, et
$A^{(r+1)}, b^{(r+1)}$ les poids et biais de la couche de sortie, inconnus de
l'attaquant. Le logit de la classe $k$ s'écrit $A_k \cdot y + b_k$, et le
réseau prédit $\arg\max_k (A_k \cdot y + b_k)$.

La méthode repose sur trois hypothèses, posées en un point de transition
entre deux classes $i$ et $j$ (un point où la classe prédite bascule de $i$ à
$j$) :

\begin{itemize}
  \item[\textbf{(H1)}] $A_i \cdot y + b_i = A_j \cdot y + b_j$ : au point de
    transition exact, les deux classes sont à égalité.
  \item[\textbf{(H2)}] $y = f_{1..r}(x)$ est calculable, c'est-à-dire que les
    couches cachées sont déjà connues. C'est une hypothèse forte et
    assumée : ce résultat n'est pas une attaque de bout en bout, c'est le
    dernier maillon d'une chaîne dont EUROCRYPT 2025 fournit les maillons
    précédents.
  \item[\textbf{(H3)}] L'oracle ne renvoie que $\arg\max_k(A_k \cdot y +
    b_k)$. Cette décision ne change pas si l'on ajoute une même constante à
    tous les scores, ni si on les multiplie tous par une même constante
    positive.
\end{itemize}

Le point essentiel est qu'aucune des trois ne mentionne l'activation des
couches cachées. (H3) ne porte que sur la règle de décision finale, un
argmax, qui est strictement la même quelle que soit la fonction qui a produit
$y$. Si (H2) tient, le problème qui reste est purement linéaire : retrouver
un classifieur linéaire à partir d'un oracle qui ne donne que son argmax.
C'est ce que montre la figure~\ref{fig:transition}.

\begin{figure}[h]
\centering
\begin{tikzpicture}[scale=1.0, every node/.style={font=\small}]
  % axes
  \draw[-{Stealth[length=2mm]}] (0,0) -- (8,0) node[right] {position le long d'une droite de sonde};
  \draw[-{Stealth[length=2mm]}] (0,-0.3) -- (0,3.3) node[above] {score};
  % two score lines
  \draw[thick, blue!60!black] (0.5,0.5) -- (7.3,2.6) node[right] {$A_i\cdot y + b_i$};
  \draw[thick, red!60!black] (0.5,2.6) -- (7.3,0.6) node[right] {$A_j\cdot y + b_j$};
  % intersection
  \coordinate (P) at (3.93,1.56);
  \fill[black] (P) circle (1.6pt);
  \draw[dashed] (P) -- (3.93,-0.25) node[below] {$p$ (point de transition)};
  % regions
  \draw[-{Stealth[length=1.8mm]}, black!50] (1.5,-0.55) -- (1.5,-0.1);
  \node[black!60] at (1.5,-0.8) {$\arg\max = j$};
  \draw[-{Stealth[length=1.8mm]}, black!50] (6.3,-0.55) -- (6.3,-0.1);
  \node[black!60] at (6.3,-0.8) {$\arg\max = i$};
\end{tikzpicture}
\caption{Au point $p$, les scores des classes $i$ et $j$ sont égaux (H1).
De part et d'autre, l'une domine l'autre : c'est tout ce que l'oracle
révèle (H3). Rien ici ne dépend de la façon dont $y$ a été calculé (H2).}
\label{fig:transition}
\end{figure}

\begin{figure}[h]
\centering
\begin{tikzpicture}[scale=0.95, every node/.style={font=\small}]
  % left: translation invariance
  \begin{scope}
    \draw[-{Stealth[length=2mm]}] (0,0) -- (3.6,0);
    \draw[-{Stealth[length=2mm]}] (0,-0.2) -- (0,2.6);
    \draw[thick, blue!60!black] (0.3,0.3) -- (3.3,1.9);
    \draw[thick, red!60!black] (0.3,1.9) -- (3.3,0.3);
    \draw[thick, blue!30!white, dashed] (0.3,0.9) -- (3.3,2.5);
    \draw[thick, red!30!white, dashed] (0.3,2.5) -- (3.3,0.9);
    \fill[black] (1.8,1.1) circle (1.4pt);
    \node at (1.8,-0.6) {translation (même $p$)};
  \end{scope}
  % right: scale with negative sign
  \begin{scope}[xshift=5.4cm]
    \draw[-{Stealth[length=2mm]}] (0,0) -- (3.6,0);
    \draw[-{Stealth[length=2mm]}] (0,-1.6) -- (0,2.6);
    \draw[thick, blue!60!black] (0.3,0.3) -- (3.3,1.9);
    \draw[thick, red!60!black] (0.3,1.9) -- (3.3,0.3);
    \draw[thick, blue!30!white, dashed] (0.3,-0.3) -- (3.3,-1.9);
    \draw[thick, red!30!white, dashed] (0.3,-1.9) -- (3.3,-0.3);
    \fill[black] (1.8,1.1) circle (1.4pt);
    \node at (1.8,-2.2) {échelle $c<0$ : $\arg\max \to \arg\min$};
  \end{scope}
\end{tikzpicture}
\caption{Les deux invariances de (H3). À gauche : décaler les deux scores de
la même constante ne change pas le gagnant. À droite : les multiplier par une
même constante négative (en pointillé) inverse l'ordre complet, donc
l'argmax devient un argmin. C'est exactement ce qui se produit pour ReLU,
GELU et SiLU dans les résultats mesurés (section~\ref{sec:resultats}).}
\label{fig:invariances}
\end{figure}

Ce constat est attendu une fois (H2) admise. L'intérêt de ce résultat n'est
pas de le deviner mais de le vérifier explicitement sur la méthode de
l'article de départ, et de traiter au passage l'ambiguïté de signe en sortie
multi-classes, que cet article ne discute que pour la sortie binaire.

\section{Protocole}

Réseau $24$-$20$-$16$-$8$-$4$ (profondeur 4, 4 classes en sortie), construit
avec la même architecture et la même graine (\texttt{seed=5}) pour les
quatre activations testées. L'attaquant dispose de l'oracle hard-label et
suppose les couches $1..r$ déjà extraites (hypothèse H2, posée ici
explicitement). Il collecte des points de transition, résout le système
linéaire qui en découle, puis corrige l'ambiguïté de signe global avec une
seule requête supplémentaire.

\subsection{Compter les degrés de liberté}

La couche de sortie compte $4 \times (8+1) = 36$ coefficients
($d_r = 8$ neurones cachés, $d_{r+1} = 4$ classes, $+1$ pour chaque biais).
L'invariance par translation (H3) en rend $8+1=9$ inobservables (on peut
décaler les 8 colonnes de $A$ et le biais sans rien changer à l'argmax), et
l'échelle globale un de plus. Le rang maximal théorique du système est donc

\[
  \text{rang}_{\max} = \underbrace{36}_{\text{inconnues}} -
  \underbrace{\big(\underbrace{9}_{\text{translation}} +
  \underbrace{1}_{\text{échelle}}\big)}_{\text{degrés de liberté, } d_r+2} = 26.
\]

La figure~\ref{fig:kernel} rend ce calcul concret sur la structure réelle du
vecteur d'inconnues $\theta$.

\begin{figure}[h]
\centering
\begin{tikzpicture}[scale=0.62, every node/.style={font=\scriptsize}]
  \def\cw{0.62} % largeur cellule
  \def\ch{0.62} % hauteur cellule
  \foreach \row [count=\r from 0] in {0,1,2,3} {
    \foreach \col in {0,...,8} {
      \pgfmathsetmacro{\x}{\col*\cw}
      \pgfmathsetmacro{\y}{-\row*\ch}
      \ifnum\row=0
        \fill[black!25] (\x,\y) rectangle ++(\cw,-\ch);
      \else
        \ifnum\col=0
          \ifnum\row=1
            \fill[blue!25] (\x,\y) rectangle ++(\cw,-\ch);
          \else
            \fill[white] (\x,\y) rectangle ++(\cw,-\ch);
          \fi
        \else
          \fill[white] (\x,\y) rectangle ++(\cw,-\ch);
        \fi
      \fi
      \draw[black!50] (\x,\y) rectangle ++(\cw,-\ch);
    }
    \node[anchor=east] at (-0.15,-\r*\ch-\ch/2) {classe \r};
  }
  \node[anchor=south] at (4*0.62,0.3) {\normalsize bloc $k$ : $A_k$ (8 valeurs) puis $b_k$ (1 valeur), répété pour $k=0..3$};
  \draw[decorate, decoration={brace, amplitude=4pt}, thick]
    (-0.3,0) -- (-0.3,-0.62) node[midway, left=6pt] {};
  \node[anchor=east, black!50] at (-1.6,-0.31) {fixé à $0$ \\ (9 val.)};
  \node[anchor=east, blue!60!black] at (-1.6,-0.93) {fixé à $1$};
\end{tikzpicture}
\caption{Les 36 inconnues de $\theta$, rangées en 4 blocs de 9 (8 coefficients
de $A$ + 1 biais, par classe). Le bloc entier de la classe 0 (gris) est fixé
à zéro : c'est le choix de représentant pour les 9 degrés de liberté de
translation. Un coefficient supplémentaire (bleu, $A_{1,1}=1$) fixe
l'échelle. Les 26 cellules blanches restantes sont les inconnues réellement
résolues, exactement le rang maximal calculé plus haut.}
\label{fig:kernel}
\end{figure}

L'« écart max » utilisé plus loin est l'erreur maximale $\max |\hat A -
A_{\text{vrai}}|$ entre les coefficients retrouvés et les vrais, une fois ces
derniers ramenés à la même normalisation (bloc 0 à zéro, $A_{1,1}=1$). Il ne
faut jamais comparer $\hat A$ directement à $A$ : les deux vivent dans la
même classe d'équivalence, mais pas au même représentant.

\section{Déroulement pas à pas, sur un cas réel (ReLU)}
\label{sec:pasapas}

Cette section suit l'exécution réelle du script sur l'activation ReLU,
graine 5, étape par étape, avec les chiffres effectivement produits. Les
trois autres activations suivent exactement le même déroulement ; seuls les
résultats numériques diffèrent, rassemblés section~\ref{sec:resultats}.

\subsection*{Étape 0 : le réseau cible et ce que l'attaquant ne voit pas}

Le réseau cible a l'architecture $24$-$20$-$16$-$8$-$4$, exactitude
$100\,\%$ sur les données d'entraînement. Sa vraie couche de sortie, que
l'attaquant ne voit jamais, est

\[
A^{(r+1)} = \begin{pmatrix}
 0{,}4161 & -0{,}3239 & -0{,}4117 & 0{,}5513 & 1{,}5229 & 0{,}3504 & 0{,}5789 & -0{,}2251 \\
-0{,}0752 & -0{,}7446 & 0{,}5031 & -0{,}1918 & -0{,}0407 & 0{,}1600 & 0{,}3785 & 0{,}7619 \\
 1{,}2533 & -0{,}3054 & -0{,}5891 & 0{,}0522 & -0{,}3394 & 0{,}5158 & -0{,}5281 & 0{,}0588 \\
 0{,}0010 & 0{,}8618 & 0{,}5219 & 0{,}8480 & 0{,}0627 & -0{,}9511 & -0{,}2533 & 0{,}2501
\end{pmatrix}, \quad
b^{(r+1)} = \begin{pmatrix} 0{,}1928 \\ -0{,}2417 \\ -0{,}1632 \\ 0{,}2121 \end{pmatrix}.
\]

\subsection*{Étape 1 et 2 : l'oracle, et la recherche d'un point de transition}

L'attaquant ne voit que $\arg\max_k f(x)_k$, jamais les logits. Deux points
$x_a$, $x_b$ sont tirés, menant respectivement aux classes $1$ et $3$
(distance initiale $\|x_b - x_a\| = 12{,}3137$). Une dichotomie resserre
l'intervalle $[x_a, x_b]$ jusqu'à un point de transition $p$ :

\begin{center}
\begin{tabular}{rl}
\toprule
itération $k$ & largeur $\|hi - lo\|$ \\
\midrule
1   & $6{,}157 \times 10^{0}$ \\
10  & $1{,}203 \times 10^{-2}$ \\
20  & $1{,}174 \times 10^{-5}$ \\
40  & $1{,}120 \times 10^{-11}$ \\
55  & $3{,}988 \times 10^{-16}$ \\
100 & $3{,}554 \times 10^{-16}$ \\
\bottomrule
\end{tabular}
\end{center}

La figure~\ref{fig:bisection} illustre ce resserrement. Le plancher atteint
vers $k=55$, environ $2{,}2 \times 10^{-16}$, est l'epsilon machine du
flottant double ($2^{-52}$) : au-delà d'une cinquantaine d'itérations,
chaque requête supplémentaire est gaspillée, la précision ne peut plus
s'améliorer.

\begin{figure}[h]
\centering
\begin{tikzpicture}[scale=1.0, every node/.style={font=\small}]
  \draw[thick] (0,0) -- (10,0);
  \node[below] at (0,0) {$x_a$ (classe 1)};
  \node[below] at (10,0) {$x_b$ (classe 3)};
  \fill (0,0) circle (1.5pt);
  \fill (10,0) circle (1.5pt);
  % k=1
  \draw[blue!70!black, thick] (0,0.35) -- (5,0.35);
  \node[blue!70!black, left] at (-0.1,0.35) {$k=1$};
  % k=5
  \draw[blue!50!black, thick] (3.2,0.75) -- (4.6,0.75);
  \node[blue!50!black, left] at (3.0,0.75) {$k=5$};
  % k=10
  \draw[blue!30!black, thick] (3.85,1.15) -- (4.05,1.15);
  \node[blue!30!black, left] at (3.65,1.15) {$k=10$};
  % p
  \draw[dashed] (3.93,0) -- (3.93,1.5);
  \fill[red] (3.93,0) circle (1.6pt);
  \node[red, above] at (3.93,1.55) {$p$ : $\|hi-lo\| \to \varepsilon_{\text{machine}}$};
\end{tikzpicture}
\caption{Chaque itération de dichotomie divise l'intervalle $[lo,hi]$ par
deux. Après une cinquantaine d'itérations, l'intervalle est plus étroit que
ce que peut distinguer un flottant double : $p$ est localisé à la précision
machine.}
\label{fig:bisection}
\end{figure}

\subsection*{Étape 3 : identifier le couple de classes}

Deux sondes de part et d'autre de $p$, à $\varepsilon = 10^{-7}$, renvoient
les classes $1$ et $3$ : c'est le couple $(i,j)$ qui se rencontre en ce
point. Si les deux sondes renvoyaient la même classe, le point serait
rejeté : la dichotomie aurait convergé près d'un coin où trois classes se
rencontrent, ou la précision ne permettrait plus de séparer.

\subsection*{Étape 4 : construire une ligne du système}

En $p$, la sortie de la dernière couche cachée vaut (connue par H2) :

\[
y = (0 ;\ 3{,}6632 ;\ 11{,}4860 ;\ 0 ;\ 0 ;\ 0{,}2082 ;\ 1{,}8252 ;\ 10{,}0999).
\]

L'équation $A_1 \cdot y + b_1 = A_3 \cdot y + b_3$ devient une ligne du
système : le bloc de la classe 1 reçoit $(y, 1)$, le bloc de la classe 3
reçoit $-(y,1)$, les blocs des classes 0 et 2 restent nuls. Une vérification
systématique confirme que la vraie couche de sortie satisfait bien cette
ligne : $\text{ligne} \cdot \theta_{\text{vrai}} = 2{,}387 \times 10^{-15}$,
proche de zéro à la précision machine près. Cette vérification est faite à
chaque exécution : elle attrape une erreur d'indice ou de signe en une
ligne, avant qu'elle ne coûte une journée de débogage.

\subsection*{Étape 5 : collecter assez de points}

L'étape 2 à 4 est répétée jusqu'à réunir $6 \times 26 = 156$ points de
transition (203 tirages tentés, 0 rejeté), formant un système de $156$
équations pour $36$ inconnues, assemblé en $0{,}9$ seconde. Les 6 couples de
classes possibles apparaissent tous, avec une fréquence qui varie de 17 à 40
occurrences : un couple absent ferait perdre du rang au système.

\subsection*{Étape 6 : vérifier le rang et le noyau}

Le rang du système atteint $26$, le maximum théorique, et la dimension du
noyau vaut $36-26=10$, exactement $d_r+2$. Chacun des 10 générateurs
attendus (une translation par colonne de $A$, une translation du biais, une
mise à l'échelle) est bien dans le noyau :

\begin{center}
\begin{tabular}{lr}
\toprule
générateur & $\|\text{système} \cdot v\|$ \\
\midrule
translation colonne 0 à 7 de $A$ (8 générateurs) & $0{,}00 \times 10^{0}$ \\
translation du biais $b$ & $0{,}00 \times 10^{0}$ \\
mise à l'échelle de $(A,b)$ & $5{,}82 \times 10^{-14}$ \\
\bottomrule
\end{tabular}
\end{center}

Les 9 premiers sont des identités structurelles, exactement nuls. Le
dernier n'est pas rigoureusement nul parce que c'est $\theta_{\text{vrai}}$
lui-même : les lignes ne sont satisfaites qu'à la précision des points de
transition, pas exactement.

\subsection*{Étape 7 : fixer 10 variables, résoudre le reste}

On fixe le bloc entier de la classe 0 à zéro (9 valeurs) et $A_{1,1}$ à $1$
(1 valeur), exactement les $10$ degrés de liberté identifiés. Le système
réduit, de rang plein $26/26$, se résout par moindres carrés en $1{,}3$~ms,
avec un résidu maximal de $1{,}036 \times 10^{-13}$. Comparée à la vérité
terrain normalisée de la même façon, l'erreur maximale sur les coefficients
retrouvés est

\[
\max |\hat A - A_{\text{vrai, normalisé}}| = 4{,}297 \times 10^{-14},
\]

de l'ordre de l'epsilon machine. Rien n'indique qu'une activation lisse ou
une activation en coude change quoi que ce soit à cette précision ; la
section~\ref{sec:resultats} le confirme sur les trois autres activations.

\subsection*{Étape 8 : l'ambiguïté de signe}

En imposant $A_{1,1}=1$, l'attaque fixe implicitement un facteur d'échelle
$c = 1 / A_{1,1,\text{vrai}}$. Ici, $A_{1,1,\text{vrai}} = -0{,}491228$,
donc $c = -2{,}035713$, négatif. Comme le montre la
figure~\ref{fig:invariances}, un facteur d'échelle négatif inverse l'ordre
complet des scores : l'argmax du réseau reconstruit devient un argmin.

L'effet est spectaculaire avant correction : $0$ accord sur $5\,000$ entrées
de test, soit $0{,}0\,\%$, très en dessous du hasard à $25\,\%$ pour 4
classes. Un score aussi inférieur au hasard n'est jamais un échec de la
méthode : c'est la signature exacte d'une convention de signe inversée. La
correction ne coûte qu'une seule requête supplémentaire, choisie à part des
points de transition : elle compare la classe prédite par $\hat A, \hat b$ à
la classe réelle sur une entrée quelconque, et inverse le signe de $\hat A,
\hat b$ si elles diffèrent. La figure~\ref{fig:signe} illustre ce
mécanisme.

\begin{figure}[h]
\centering
\begin{tikzpicture}[scale=1.0, every node/.style={font=\small}]
  % original order
  \node[anchor=east] at (-0.3,1.5) {scores $u$ :};
  \foreach \x/\v/\lbl in {0/0.3/classe 2, 2/0.9/classe 0, 4/1.5/classe 1, 6/0.6/classe 3} {
    \fill[blue!60!black] (\x,1.5*\v) circle (2pt);
    \node[below] at (\x,0) {\lbl};
  }
  \draw[-{Stealth[length=1.8mm]}, thick] (4,2.3) -- (4,2.6) node[above] {$\arg\max(u) = $ classe 1};
  % arrow down to c*u
  \draw[-{Stealth[length=2mm]}, black!50] (3,-0.6) -- (3,-1.3) node[midway, right] {$\times\, c,\ c<0$};
  % flipped order
  \node[anchor=east] at (-0.3,-2.5) {scores $c\cdot u$ :};
  \foreach \x/\v/\lbl in {0/-0.3/classe 2, 2/-0.9/classe 0, 4/-1.5/classe 1, 6/-0.6/classe 3} {
    \fill[red!60!black] (\x,-2.5-1.5*(-\v)) circle (2pt);
  }
  \draw[-{Stealth[length=1.8mm]}, thick] (0,-4.0) -- (0,-4.3) node[below] {$\arg\max(c\cdot u) = $ classe 2 $= \arg\min(u)$};
\end{tikzpicture}
\caption{Avec un facteur d'échelle positif imposé arbitrairement
($A_{1,1}=1$), rien ne garantit que le vrai facteur soit positif. S'il est
négatif, comme pour ReLU, GELU et SiLU dans ces résultats, le réseau
reconstruit classe à l'envers : l'argmax devient un argmin. Une seule
requête de sonde suffit à le détecter et à inverser les signes.}
\label{fig:signe}
\end{figure}

\subsection*{Étape 9 : validation de l'équivalence fonctionnelle}

Après correction, l'accord est parfait sur deux jeux de test indépendants :
$5\,000/5\,000$ sur des entrées uniformes, $5\,000/5\,000$ sur des entrées
gaussiennes. Tester sur deux distributions n'est pas redondant : une
équivalence qui ne tiendrait que sur la distribution d'entraînement serait
locale, pas fonctionnelle. Les matrices $\hat A$ et $A$ restent différentes
(on a récupéré une classe d'équivalence, pas $\theta$ lui-même), mais les
arg\-max coïncident partout.

\subsection*{Étape 10 : où part le coût}

\begin{center}
\begin{tabular}{lrr}
\toprule
phase & requêtes & part \\
\midrule
recherche des points de transition & 16\,581 & $100{,}0\,\%$ \\
exploration (trouver le premier point) & 1 & $0{,}0\,\%$ \\
correction du signe & 1 & $0{,}0\,\%$ \\
\midrule
total & 16\,583 & \\
\bottomrule
\end{tabular}
\end{center}

La résolution du système linéaire lui-même ne coûte aucune requête, et
quelques millisecondes. Presque tout le coût est dans la recherche des
points de transition, et c'est strictement la même opération quelle que
soit l'activation cachée : elle ne regarde que la couche de sortie, jamais
ce qu'il y a avant. L'exécution complète prend $8{,}1$ secondes.

\begin{remarque}[Pourquoi cette étape intéresse autant l'attaque que la défense]
Côté attaque, cette recherche est trivialement parallélisable : chaque point
de transition est indépendant des autres. Côté défense, c'est précisément
la cible à viser. L'attaque a besoin d'une précision de $10^{-13}$ sur la
position du point de transition, une contrainte numériquement
extraordinairement fragile : bruiter la sortie de l'oracle, même très
légèrement, romprait cette précision et donc la méthode.
\end{remarque}

\section{Résultats mesurés sur les quatre activations}
\label{sec:resultats}

Le protocole de la section précédente est rejoué à l'identique (même
architecture, même graine $5$) pour GELU, SiLU et tanh. Les quatre
exécutions sont indépendantes, chacune journalisée intégralement.

\begin{table}[h]
\centering
\begin{tabular}{lrrrrr}
\toprule
Activation & Rang atteint & Écart max & Échelle $c$ & Accord avant & Accord après \\
\midrule
ReLU  & $26/26$ & $4{,}297 \times 10^{-14}$ & $-2{,}035713$ (négatif) & $0/5000$ & $5000/5000$ \\
GELU  & $26/26$ & $1{,}306 \times 10^{-13}$ & $-2{,}831870$ (négatif) & $0/5000$ & $5000/5000$ \\
SiLU  & $26/26$ & $2{,}132 \times 10^{-14}$ & $-4{,}920914$ (négatif) & $0/5000$ & $5000/5000$ \\
tanh  & $26/26$ & $6{,}217 \times 10^{-15}$ & $0{,}548012$ (positif)  & $5000/5000$ & $5000/5000$ \\
\bottomrule
\end{tabular}
\caption{Résultat principal. Le rang atteint et l'écart résiduel sont
identiques, à la précision machine près, pour les quatre activations. Le
signe de $c$ ne dépend pas de l'activation : c'est un hasard de tirage sur
$A_{1,1,\text{vrai}}$, pas un effet systématique (3 négatifs, 1 positif sur
4 essais).}
\end{table}

\begin{table}[h]
\centering
\begin{tabular}{lrr}
\toprule
Activation & Requêtes totales & Temps \\
\midrule
ReLU & 16\,583 & $8{,}1$ s \\
GELU & 16\,581 & $16{,}4$ s \\
SiLU & 16\,589 & $9{,}3$ s \\
tanh & 16\,577 & $8{,}4$ s \\
\bottomrule
\end{tabular}
\caption{Le nombre de requêtes est quasi identique entre les quatre
activations, ce qui confirme que la recherche des points de transition ne
dépend pas de l'activation cachée. Le temps de GELU, environ deux fois plus
long, s'explique par le coût de simulation de l'oracle (la fonction
$\mathrm{erf}$ y est réimplémentée en Python pour rester sans dépendance à
scipy), pas par l'attaque elle-même.}
\end{table}

\subsection*{Lecture}

Le rang atteint, $26$ sur $26$, le maximum théorique pour cette
architecture, et l'écart résiduel, de l'ordre de $10^{-13}$ à $10^{-15}$,
proche de la précision machine, sont identiques pour les quatre
activations. C'est la confirmation directe de ce qu'annonçaient les trois
hypothèses : aucune ne porte sur l'activation cachée, et le comportement
mesuré ne varie pas avec elle.

Pour ReLU, GELU et SiLU, le facteur d'échelle global $c$ ressort négatif :
la méthode reconstruit alors le réseau à l'envers, d'où un accord de
$0/5000$ avant correction. Ce n'est pas un échec, c'est la signature
attendue d'une inversion de convention (figure~\ref{fig:signe}), et une
seule requête supplémentaire suffit à la détecter et à la corriger. Pour
tanh, $c$ ressort positif d'emblée, donc aucune correction n'est
nécessaire. Après correction, l'accord est parfait sur deux distributions
de test, pour les quatre activations.

\section{Limites}

\begin{itemize}
  \item Le résultat suppose les couches cachées $1..r$ déjà connues
    (hypothèse H2). Ce n'est pas encore une attaque de bout en bout sur un
    réseau dont tout est inconnu.
  \item Testé sur un seul réseau jouet par activation, une seule graine. Le
    réseau CIFAR-10 réel de l'article de départ n'a pas encore été rejoué
    avec cette méthode.
  \item La correction de l'ambiguïté de signe global en sortie multi-classes
    (étape 8) va au-delà de ce que traite l'article de départ, qui ne la
    discute que pour la sortie binaire. C'est une observation empirique
    mesurée ici sur 4 exécutions indépendantes, pas encore démontrée
    formellement.
\end{itemize}

\section{Reproduire ces résultats}

\begin{center}
\begin{tabular}{ll}
\toprule
Activation & Journal (dans \texttt{logs/}) \\
\midrule
ReLU & \texttt{couche\_sortie\_relu.txt} \\
GELU & \texttt{couche\_sortie\_gelu.txt} \\
SiLU & \texttt{couche\_sortie\_silu.txt} \\
tanh & \texttt{couche\_sortie\_tanh.txt} \\
\bottomrule
\end{tabular}
\end{center}

Pour rejouer, depuis \texttt{travail\_aout\_2026/07-08-2026/} :

\begin{quote}
\verb|python attaque_couche_sortie.py --activation <nom> --seed 5 --sans-diagnostic|
\end{quote}

Le script fixe sa graine : les chiffres de ce document se reproduisent à
l'identique à chaque exécution.

\section{Prochaines étapes sur ce résultat}

\begin{enumerate}
  \item Démontrer formellement la correction de signe en sortie
    multi-classes (étape 8), au lieu de la mesurer seulement
    empiriquement.
  \item Rejouer l'attaque sur plusieurs graines, pour vérifier que les
    chiffres de la section~\ref{sec:resultats} ne sont pas un effet de
    tirage favorable.
  \item Rejouer l'attaque sur le réseau CIFAR-10 de l'article de départ, pour
    passer d'un réseau jouet à l'échelle réelle utilisée dans la
    littérature.
\end{enumerate}

\end{document}
